\section{Two-dimensional muon cross section}
\label{sec:part-2d}
\label{sec:results}

\subsection{Unfolded cross section and comparison with MINERvA}
\label{sec:muon-cross-section-comparison}
\label{sec:results-xsec}

We use the two-dimensional refined OmniFold estimator of
Sec.~\ref{sec:estimators-and-extraction} to extract the flux-averaged inclusive
charged-current $\nu_\mu$ cross section per nucleon on hydrocarbon,
$d^2\sigma/(dp_T\,dp_\parallel)$. The 205 reported bins match the published
MINERvA phase space and binning~\cite{MINERvA:2021owq}. Figure~\ref{fig:muon-projections}
shows the two momentum projections, obtained by integrating the unfolded result
over the other momentum component. Both projections follow the published spectra
through the peak and the falling tails. In every $p_T$ slice and every
$p_\parallel$ slice, the unfolded spectrum is closer to the published cross section
than to the prior. \phantomsection\label{sec:results-model}The published GENIE
2.12.6 prediction with MINERvA Tune v1 gives about $\fpeval{100*0.91}\%$ of the
published integrated cross section, with the deficit
concentrated near the peak of the $p_T$ spectrum. The two unfolded results show
the same excess over that prediction in the peak region. In the lowest $p_T$
bin, the OmniFold projection falls below the prior, and the ratio rises
above unity toward higher $p_T$.

\begin{figure}[H]
  \centering
  \includegraphics[width=0.49\textwidth]{MEFHC_5iter_xsec_proj_pt}
  \includegraphics[width=0.49\textwidth]{MEFHC_5iter_xsec_proj_pz}
  \caption{Muon transverse (left) and longitudinal (right) momentum projections
  of the unfolded cross section, compared with the published MINERvA
  result~\cite{MINERvA:2021owq}. The solid and dashed histograms show GENIE
  2.12.6 with MINERvA Tune v1 for the unfolding prior and the published
  prediction, respectively. The lower panels give the ratio of the OmniFold
  result to its prior. Horizontal bars span the bins. The vertical bars on
  the OmniFold points show finite weighted-sample statistics only.}
  \label{fig:muon-projections}
  \label{fig:xsec}
\end{figure}

Integrating over the reported bins gives $\num{\sigTwoD}\,\mathrm{cm}^2$
per nucleon, compared with $\num{\sigTwoDpaper}\,\mathrm{cm}^2$ per nucleon
for MINERvA, a ratio of \ratioTot. This integration weights the differential
cross section in each bin by its area. The fractional uncertainty from the
covariance of the OmniFold result is $\uqIntegral\%$ on the integral, and the
corresponding published uncertainty is $\uqPaperIntegral\%$. Agreement of the integrals
does not test the shape, because excesses and deficits in different regions
can cancel.

\phantomsection\label{sec:results-compare}
The median binwise ratio to MINERvA is \medianBinRatio, and $\binsFive\%$, $\binsTen\%$, and
$\binsTwenty\%$ of the bins agree within $5\%$, $10\%$, and $20\%$, respectively.
Figure~\ref{fig:muon-standardized-differences} shows the differences in units
of the published bin uncertainties. The distribution has a mean of \pullMean{} and a
standard deviation of \pullRMS. Let $\Delta$ be the vector of differences between
the differential cross sections and $V_{\mathrm{pub}}$ the published covariance.
Using only the diagonal of $V_{\mathrm{pub}}$ gives a quadratic distance per
bin of 0.29. Using the full matrix gives
$\Delta^{\mathsf T}V_{\mathrm{pub}}^{+}\Delta/205=\chiPaper$, where the superscript
$+$ denotes the pseudoinverse. Thus the covariance correlations give
substantially greater weight to the observed pattern of differences than the
diagonal comparison suggests. The published uncertainty of a linear
combination of bins can be small even when the uncertainties of its
constituent bins are comparatively large. This makes the quadratic distance
sensitive to differences that are small on the scale of the individual bins.

\begin{figure}[H]
  \centering
  \includegraphics[width=\textwidth]{MEFHC_5iter_pull_full}
  \caption{Difference between the OmniFold and published cross sections in each
  reported bin, divided by the published standard deviation for that bin. The
  left panel shows the differences on the muon momentum grid, and the right
  panel shows their distribution. The map and histogram use the diagonal
  uncertainties. The value of $\chi^2$ per bin in the inset uses the full
  published covariance.}
  \label{fig:muon-standardized-differences}
  \label{fig:pull}
\end{figure}

\phantomsection\label{sec:results-tension}
The distance is dominated by shape. Using the published covariance, the
normalization difference gives $\chi^2_{\mathrm{norm}}=0.04$, whereas comparison
of the unit-area shapes with the propagated covariance gives
$\chi^2_{\mathrm{shape}}\simeq465$. The shape difference contributes across many
covariance eigenmodes, including both small- and moderate-variance directions.
A comparison of separate run-period unfolds finds that the observed run-period
step cannot account for the distance in the best-determined covariance
directions at the sensitivity of that test.
Appendix~\ref{app:comparison-closure-coverage} defines the comparison statistics, and
Appendix~\ref{app:additional-2d-comparisons} gives the
additional comparisons.

\phantomsection\label{sec:results-benchmarks}
A binned IBU of the same inputs at ten iterations, the setting of the published
measurement~\cite{MINERvA:2021owq}, differs from the two-dimensional OmniFold
result beyond paired-replica fluctuations ($T=\fvval{c2c1-T}$, compared with a null
95th percentile of \fvval{c2c1-q95}). Relative to the published result and using its
covariance, ten-iteration IBU and OmniFold give descriptive $\chi^2$ per bin of
\fvval{c2c1-dibu} and \fvval{c2c1-der}, respectively. The origin of the correlated shape difference
remains undetermined.

The comparison with the publication is descriptive. The two results share data
and flux, detector, and interaction-model uncertainties, and their
cross-covariance is unavailable. With $K$ the cross-covariance between the
OmniFold and published results, the covariance of their difference is
\begin{equation}
  \operatorname{Cov}(\Delta)=Q_R+V_{\mathrm{pub}}-K-K^{\mathsf T} .
  \label{eq:muon-difference-covariance}
\end{equation}
Using the simple sum $Q_R+V_{\mathrm{pub}}$ gives a distance per bin of \chiCombined.
This sum omits the cross terms and counts shared uncertainty contributions
twice. The published covariance is also ill-conditioned, because the variance differs
by many orders of magnitude across its eigenmodes. These distances therefore
carry no goodness-of-fit probability. The tree implementation and the
refined background subtraction were adopted with the published comparison
available (Appendix~\ref{app:choices-after-results}). Together these conditions make the publication
comparison a benchmark of the central values with shared inputs.

\subsection{Uncertainty and covariance}
\label{sec:muon-cross-section-covariance}
\label{sec:systematics}
\label{sec:uq-combine}

The reported $205\times205$ covariance $Q_R$ describes the differential cross
section obtained with the same estimator as the central value. It is the sum
of the systematic, statistical, training-seed, and run-period contributions,
\begin{equation}
  Q_R=C_{\mathrm{syst}}+C_{\mathrm{stat}}+C_{\mathrm{seed}}+C_{\mathrm{run}} .
  \label{eq:muon-covariance-block-sum}
\end{equation}
The sum assumes zero cross-covariances between blocks and a linear response to
simultaneous perturbations. The common data and unfolding can couple the effects
of different sources of variation. For example, changes in the training sample
can affect the dependence of the result on the training seed. The reported sum
has no separate block for interactions of seed variation with the bootstrap or
systematic perturbations. Appendix~\ref{app:covariance-construction} gives the construction of the
ensembles, the treatment of the flux normalization,
and the assumptions used to combine their covariances.

\phantomsection\label{sec:results-uq}
\begin{table}[H]
  \caption{Uncertainty budget for the 205 reported muon bins. For each
  covariance, the percentage is the median of the binwise fractional standard
  deviations. The square-root trace gives the diagonal uncertainties in
  quadrature, in units of $10^{-39}\,\mathrm{cm}^2$ per nucleon per
  $(\mathrm{GeV}/c)^2$. The published entry uses the published central values
  and covariance~\cite{MINERvA:2021owq}.}
  \label{tab:muon-uncertainty-budget}
  \label{tab:uq}
  \centering
  \begin{tabular}{@{}lcc@{}}
    \toprule
    Source & Median (\%) & $\sqrt{\operatorname{tr}C}$ \\
    \midrule
    Systematics and normalization & \num[round-mode=places,round-precision=2]{5.926}
      & \fpeval{2.686e-39 * 1e39} \\
    Statistical bootstrap & 0.676 & \fpeval{1.886e-40 * 1e39} \\
    Training-seed variation & 0.048 & \fpeval{1.351e-41 * 1e39} \\
    Run-period differences & \num[round-mode=places,round-precision=2]{2.087}
      & \fpeval{8.314e-40 * 1e39} \\
    Combined ($Q_R$) & \uqMedian & \fpeval{2.818e-39 * 1e39} \\
    Published total & \uqPaper & \fpeval{2.676e-39 * 1e39} \\
    \bottomrule
  \end{tabular}
\end{table}

\phantomsection\label{sec:uq-universes}
The systematic block in Table~\ref{tab:muon-uncertainty-budget} includes
43 bands of simulated flux,
detector, and interaction-model variations and a fully correlated $1.4\%$
term for the target-nucleon normalization. The band sum omits a duplicate
nonresonant-pion variation.
The flux is the largest individual band, with a median fractional standard
deviation of $\fluxBand\%$, compared with $4.01\%$ in the publication. The next
largest are the MINOS muon energy scale ($2.36\%$), the MINOS muon-match
efficiency ($1.50\%$), and the MINERvA muon energy scale ($1.44\%$).
The flux prediction affects both the distribution of simulated events and the
conversion from unfolded yields to cross sections. The MINOS energy-scale
variations already incorporate correlated flux effects through their
normalization, as discussed in Appendix~\ref{app:covariance-construction}.
Under the MINERvA Analysis Toolkit (MAT) convention, each systematic band is
mean-centered and uses population normalization ($1/N$ for equal weights),
with constraint weights for the flux universes, so the band excludes the
offset of its mean from the central value.
In the systematic re-unfoldings, each missed signal event carries the truth
weight of the corresponding universe.

\phantomsection\label{sec:uq-boot}
The statistical block propagates finite data and simulation statistics through
the full unfolding using 300 Poisson bootstrap replicas at the central training
seed. Its square-root trace is 0.407 times that of the published statistical
covariance, and the median ratio of the binwise statistical standard deviations
is 0.425. The bootstrap covariance omits the replica mean's offset from the
central value, a median of 0.25 statistical standard deviations across bins
that leaves the bias relative to truth undetermined.

\phantomsection\label{sec:uq-ml}
The training-seed block measures the variation attributable to
randomized training on matched resamples and has the smallest median
contribution. It is estimated from five seed settings, which limits the
precision of the estimated variance.

The rank-two run-period block represents the spread of three period groups,
including an unexplained $(\fvval{rp-pb-unf}\pm\fvval{rp-pb-unf-se})\%$ step at $p_\parallel=8\,\mathrm{GeV}/c$
in the double ratio of unfolded to simulated cross sections for the middle and
earliest groups, and raises the uncertainty on the integral by $9.9\%$.

For the projections, the covariance is $P Q_R P^{\mathsf T}$, where $P$
integrates over the removed momentum component in each projected bin. The
variance of the integral is $a^{\mathsf T}Q_R a$, where $a$ is the vector of bin
areas. These quantities use the full covariance, including correlations between
bins that enter the same sum. A fully correlated normalization term, for example,
retains its fractional size on both projections and on the integral. The smaller
median bin uncertainty compared with MINERvA can therefore coexist with a larger
fractional uncertainty on the integral.

Appendix~\ref{app:covariance-inversion} discusses how small eigenvalues affect the inversion of $Q_R$.

\phantomsection\label{sec:uq-2d-scope}
We report the prior-shape dependence separately as $C_{\mathrm{model}}$, the
sum of outer products of shifts obtained by reweighting the prior to three
alternative generator shapes and repeating the unfolding. Its fractional standard
deviation varies across the grid, with a median of $\cmodelMedian\%$, an 84th percentile
of $7.10\%$, and a maximum of $21.10\%$. The median for
$Q_R+C_{\mathrm{model}}$ is $\qrCmodelMedian\%$. Steps in the reweighting ratios affect
the size of this term, with about $\fpeval{round(59.4,0)}\%$ of its trace coming from 52 bins
adjacent to such steps. Its effect on the integral is $\cmodelIntegral\%$. The
reweightings vary only the shape of the prior and leave its normalization
dependence untested. The term measures sensitivity to the chosen shapes,
with no calibrated bound on estimator bias.

$Q_R$ is conditional on the chosen learner, regularization, and parametrized
detector response. It omits variation across learner families and
regularization choices, response discrepancies outside that parametrization,
and possible bias from pooling inconsistent run periods. The construction
propagates the specified sources of variation with their correlations across
bins. The coverage studies in Sec.~\ref{sec:muon-cross-section-validation}
leave uncertainties for the reported bins, projections, and integral without
validated coverage.

\subsection{Closure tests and uncertainty validation}
\label{sec:muon-cross-section-validation}
\label{sec:validation}

We use closure tests to measure how well the unfolding recovers specified
changes to a simulated truth. Coverage tests instead measure how often
intervals or joint regions from repeated pseudo-experiments contain the target
truth. The designs and scoring rules are given in Appendix~\ref{app:comparison-closure-coverage}.

The injected-feature tests use five-iteration OmniFold with the purity-based
estimator of Table~\ref{tab:estimators} and simulation from one run period.
Reweighting both truth and reconstructed events by a Gaussian in true $p_T$
or by a tilt in true $p_\parallel$ gives median absolute fractional residuals
of $0.046\%$ and $0.013\%$, respectively, within the descriptive $1.5\%$
criterion. Scored with the feature-recovery test of
Ref.~\cite{Canelli:2025guide}, a $17.2\%$ Gaussian truth feature leaves a
$1.69\%$ residual, a ratio of 0.098 that meets the guide's criterion. Reweighting the pseudo-data by a
quasielastic axial-mass variation or a flux variation gives residuals of
$0.068\%$ and $0.376\%$, respectively, within the $2\%$ criterion when
compared with the corresponding reference extrapolated from the accepted
signal. These references probe shape changes conditional on the simulated
acceptance.

To test sensitivity to a change in the effective response, we reweight
reconstructed events as a function of the resolution variable
$\Delta p_T=p_T^{\mathrm{reco}}-p_T^{\mathrm{true}}$ at fixed truth. The two
momenta are classifier inputs at their respective levels, and their eventwise
difference is outside the feature set. For reweighting amplitudes
of 0.05, 0.10, and 0.30, the median signed fractional residuals on the
reported grid are $0.57\%$, $1.14\%$, and $3.43\%$, respectively. The two
smaller amplitudes satisfy the $3\%$ criterion, and the largest exceeds it.
The residual grows approximately linearly over this scan. All of the closure
tests use the same detector simulation, and the conclusions apply to the
specific injected shapes and amplitudes. We assign no additional uncertainty
from these closure residuals.

The background treatment is tested separately with 20 pseudo-experiments at
nominal truth, each unfolded with both treatments. The two results differ in
the mean by more than one statistical standard deviation in 14 of the 205
bins. In those bins, the purity-based estimate deviates from the truth by as
much as 5.9 statistical standard deviations, while the refined estimate is
within 0.44, using the purity-based statistical block as a common scale
(Appendix~\ref{app:negative-weight-background}).

\phantomsection\label{sec:coverage}
The quoted statistical block has had no coverage test. A test of an earlier
variant of the band is described in App.~\ref{app:choices-after-results}.

\phantomsection\label{sec:coverage-tunes}
The evaluated two-dimensional covariance under-covers jointly at both held-out
tunes. We tested \lrval{cov2d-ntoys} pseudo-experiments at each of two GENIE v3 tunes,
\fvval{f1-tune-ar23} and \fvval{f1-tune-g18}, with covariance rebuilt at the
corresponding truth, including the generator-shape term and excluding the
run-period block. The pseudo-experiments draw the variations prescribed by
the procedure at each tune. The joint region uses the $\chi^2$ quantile for the
probability content of a Gaussian one-standard-deviation interval. At both
tunes, \ifnum\lrval{cov2d-ar23-joint}=0 zero\else\lrval{cov2d-ar23-joint}\fi{}
of \lrval{cov2d-ntoys} pseudo-experiments lie in that region, below the
prespecified minimum counts of \lrval{cov2d-ar23-kjoint} and \lrval{cov2d-g18-kjoint},
respectively. The mean $\chi^2$ per degree of freedom is
$\lrval{cov2d-ar23-chi2ndf}\pm\lrval{cov2d-ar23-chi2ndf-se}$ for
\lrval{cov2d-ar23-ndf} degrees of freedom at the first tune and
$\lrval{cov2d-g18-chi2ndf}\pm\lrval{cov2d-g18-chi2ndf-se}$ for
\lrval{cov2d-g18-ndf} at the second. A descriptive
check retaining only the \lrval{t2a-r1} or \lrval{t2a-r2} largest-variance directions also finds
zero joint containments at each tune.

\begingroup
% Format the source's combined ranges as inequalities and percentage points.
\def\MuonCoverageCell#1--#2, #3--#4\relax{%
  $#1<p_T<#2\,\mathrm{GeV}/c$ and $#3<p_\parallel<#4\,\mathrm{GeV}/c$}
\def\MuonCoveragePoints#1--#2\relax{\fpeval{100*(#1)}--\fpeval{100*(#2)}}
The per-bin test scores containment within one standard deviation. It flags
\lrval{cov2d-ar23-nflagged} bins at the first tune and \lrval{cov2d-g18-nflagged} at the second, with the
deficits attributed to bias in those bins. The bin with
\expanded{\noexpand\MuonCoverageCell\lrval{cov2d-cell29-range}\noexpand\relax}
is flagged at both tunes, with \lrval{cov2d-ar23-29-count} and \lrval{cov2d-g18-29-count}
containments, respectively, against a threshold of \lrval{cov2d-ar23-29-kc}. The per-bin test
resolves deficits in coverage of about
\expanded{\noexpand\MuonCoveragePoints\lrval{cov2d-res-cell}\noexpand\relax}
percentage points. The joint test scores the full residual vector and resolves a deficit of about
\fpeval{100*\lrval{cov2d-res-joint}}
percentage points. These results apply to the covariance evaluated at the
two simulated truths. The intervals on data use $Q_R$ as defined in
Eq.~\eqref{eq:muon-covariance-block-sum}, and their coverage remains unvalidated.
\par
\endgroup

The benchmark establishes that unbinned OmniFold reproduces the published
integral to about $1\%$ and recovers specified features in simulation. A
correlated shape difference remains, and the joint shape and interval coverage
lack independent validation.

% PASS2:
% Source literals retained in 3.1: GENIE 2.12.6, Tune v1, 205 bins,
% the Tune v1 rate (0.91 as 91%), the 5%, 10%, and 20% windows, 0.29,
% the shape distances 0.04 and 465, and the null 95th percentile designation.
% Source literals retained in 3.2: the six square-root traces (rescaled to
% the draft's units), the 5.926 and 2.087 medians (rounded as in the draft),
% the 0.676 and 0.048 medians, 43 bands, 1.4%, 4.01%, 2.36%, 1.50%, 1.44%,
% 300 bootstrap replicas, 0.407, 0.425 (app_statmethods.tex), 0.25, the
% step location 8 GeV/c, its 9.9% effect on the integral, and the prior-shape
% values 84th, 7.10%, 21.10%, 59.4% (rounded to 59%), and 52.
% Source literals retained in 3.3: 0.046%, 0.013%, 1.5%, 17.2%, 1.69%, 0.098,
% 0.068%, 0.376%, and 2% (part5_validation.tex); hidden-variable amplitudes
% 0.05, 0.10, 0.30, residuals 0.57%, 1.14%, 3.43%, and threshold 3%;
% 20 pseudo-experiments (app_negweight.tex), 14 of 205, 5.9, and 0.44
% (sec_summary.tex), GENIE v3, and the conclusion's approximately 1%.
% The explicitly spelled-out count in the truncated test is also spelled
% out in sec_validation.tex. Other numeric notation is mathematical or units.
% No unresolved references in this section. Original floats without a
% retained counterpart here: tab:compare, tab:pl-added,
% tab:validation-questions, tab:coverage-bands, tab:coverage-tunes,
% fig:modelmaps, fig:uqbands, fig:bootstrap, fig:eff. The slices and tension
% figures have Appendix C.1 destinations in new-labels.md. The omitted
% diagnostic subsections have no local anchors for sec:uq-2d-construction
% or sec:val-other.
% Feature-recovery stage check: closure_truth_reweight.py's default
% gauss_pt output is the same file bottom_line_test.py reads for the 2D
% recovery score. The 17.2% feature and 1.69% residual describe the
% feature bins under that score, while 0.046% is the all-bin median residual.
% IBU placement check: C.1 preserves 19g's descriptive fixed-setting
% comparison and published metric. DECISION-20261120a item 1 released
% the sentence count only and kept the earlier placement; the later
% author-approved move in inputs2/notes.md supplies the C.1 placement.
% D.4 check: the staged appendix preserves the closure-before-comparison
% sequence. The fixed-truth numbers are omitted here per the move and are
% still absent from the staged D.4 checked during this pass.
